% !TeX root = ../../Book3.tex
% 纲要标题：### 2.5 局部环与局部同态
% 纲要位置：工作记录/计划纲要.md，第 1868 行。
% 纲要细目（写作范围）：
% * 局部环 \((R,\mathfrak m,\kappa)\) 的单位判据、极大理想与剩余域；局部性不蕴含 Noether 性
% * 局部同态及诱导的剩余域映射；有限型仿射代数的点处局部环
% * \(k[x_1,\ldots,x_n]_{(x_1,\ldots,x_n)}\)、离散赋值环及 Artin 局部环的对照
% * 先区分局部化、Hensel 化与完备化的任务；后续第13、23章分别建立相应理论
% ---

\section{局部环与局部同态}
\label{b3:sec:0205}

在 \(\mathfrak p\) 处局部化后，所有不属于 \(\mathfrak p\) 的元素都成为单位，留下的非单位恰好组成一个极大理想。这种情形值得独立命名，因为接下来的生成元判据和维数理论都将在这样的环上表述。局部性只说明单位与非单位的分界，不包含任何隐含的有限性假设。

\subsection{局部环的单位、极大理想与剩余域}
\label{b3:subsec:0205:01}

\begin{definition}\label{b3:def:local-ring}
设 \(R\) 为非零交换含幺环。若 \(R\) 恰有一个极大理想 \(\mathfrak m\)，则称 \(R\) 为\term[jubu-huan]{局部环}[local ring]。写作 \((R,\mathfrak m)\)，或连同剩余域 \(\kappa=R/\mathfrak m\) 写作 \((R,\mathfrak m,\kappa)\)。
\end{definition}

\begin{proposition}\label{b3:prop:local-ring-units}
设 \(R\) 为非零交换含幺环。下列条件等价：
\begin{enumerate}
\item \(R\) 为局部环；
\item 全部非单位组成一个理想；
\item 存在真理想 \(\mathfrak m\)，使 \(R\setminus\mathfrak m=R^\times\)。
\end{enumerate}
在这些条件下，该非单位理想就是唯一的极大理想；特别地，\(a\in\mathfrak m\) 时 \(1-a\) 是单位。
\end{proposition}
\begin{proof}
一个元素 \(a\) 非单位，当且仅当主理想 \((a)\) 为真理想，也就当且仅当 \(a\) 属于某个极大理想。若极大理想唯一，全部非单位恰为它，得到第二项与第三项。反过来，若非单位集是理想，则它包含每个真理想中的元素，又不含 \(1\)，所以是唯一的极大理想。第三项也给出同样结论。最后，若 \(a,1-a\) 都属于 \(\mathfrak m\)，则 \(1\in\mathfrak m\)，矛盾。
\end{proof}

局部环不必是 Noether 环。取
\[
R=k[x_1,x_2,\ldots]_{(x_1,x_2,\ldots)}.
\]
这是由\cref{b3:prop:localized-prime-ideals}得到的局部环，但理想链
\((x_1)\subsetneq(x_1,x_2)\subsetneq\cdots\) 仍严格上升。若 \(x_{n+1}\) 属于前一项，清分母后有
\(s x_{n+1}\in(x_1,\ldots,x_n)\)，其中 \(s\) 的常数项非零；令前 \(n\) 个变量为零，在剩余变量的多项式整环中得到一个非零多项式乘 \(x_{n+1}\) 为零，矛盾。

\subsection{局部同态与点处局部环}
\label{b3:subsec:0205:02}

\begin{definition}\label{b3:def:local-homomorphism}
设 \((R,\mathfrak m)\) 与 \((A,\mathfrak n)\) 为局部环，\(\varphi:R\to A\) 为保单位环同态。若 \(\varphi(\mathfrak m)\subseteq\mathfrak n\)，则称 \(\varphi\) 为\term[jubu-tongtai]{局部同态}[local homomorphism]。
\end{definition}

任意含幺同态都把单位送到单位，所以总有
\(\varphi^{-1}(\mathfrak n)\subseteq\mathfrak m\)。因此局部同态的条件等价于
\(\varphi^{-1}(\mathfrak n)=\mathfrak m\)，也等价于非单位不会变成单位。它诱导一个含幺的剩余域同态
\[
R/\mathfrak m\longrightarrow A/\mathfrak n,
\]
该映射为单射，因为域到非零环的含幺同态核为零。

\begin{proposition}\label{b3:prop:induced-local-map}
设 \(R,A\) 为交换含幺环，\(\varphi:R\rightarrow A\) 为保单位环同态，\(\mathfrak q\in\operatorname{Spec}A\)，并令
\(\mathfrak p=\varphi^{-1}(\mathfrak q)\)，则有自然的局部同态
\[
R_{\mathfrak p}\longrightarrow A_{\mathfrak q},
\qquad r/s\longmapsto\varphi(r)/\varphi(s).
\]
它诱导剩余域的嵌入 \(\kappa(\mathfrak p)\rightarrow\kappa(\mathfrak q)\)。
\end{proposition}
\begin{proof}
\(s\notin\mathfrak p\) 意味着 \(\varphi(s)\notin\mathfrak q\)，因而目标分母合法；局部化的泛性质给出同态。若 \(r\in\mathfrak p\)，其像落在 \(\mathfrak qA_{\mathfrak q}\)，故该同态为局部同态。再取剩余域即得所述嵌入。
\end{proof}

例如有限型仿射代数 \(A=k[x_1,\ldots,x_n]/I\) 在素理想 \(\mathfrak p\) 处的局部环为 \(A_{\mathfrak p}\)。对一个 \(k\) 有理点 \(a=(a_1,\ldots,a_n)\)，即所有 \(f\in I\) 在 \(a\) 处取零，代入映射的核 \(\mathfrak m_a\) 为极大理想，剩余域为 \(k\)。局部环中的元素可写成 \(f/g\)，要求 \(g(a)\ne0\)。这个描述表示在该点附近允许分母不消失，而不要求 \(g\) 在所有点上都非零。

\ProseParagraphBreak
局部环之间的任意含幺同态未必局部。例如 \(\mathbb Z_{(p)}\hookrightarrow\mathbb Q\) 把非单位 \(p\) 变为单位，其极大理想不能映入 \(\mathbb Q\) 的极大理想 \(0\)。所以使用剩余域映射前必须检查局部条件。

\subsection{多项式局部环、离散赋值环与 Artin 局部环}
\label{b3:subsec:0205:03}

在
\[
R=k[x_1,\ldots,x_n]_{(x_1,\ldots,x_n)}
\]
中，唯一极大理想由 \(x_1,\ldots,x_n\) 生成，剩余域为 \(k\)。当 \(n=1\) 时，每个非零元素都唯一写成 \(x^r u\)，其中 \(r\ge0\)、\(u\) 为单位。因为多项式分子可提取唯一的 \(x\) 次幂，而局部化的分母常数项非零。非零理想中取这样的最小 \(r\)，便知该理想等于 \((x^r)\)。

\begin{definition}\label{b3:def:discrete-valuation-example}
设 \(K\) 为域，满射 \(v:K^\times\rightarrow\mathbb Z\) 满足
\[
\begin{aligned}
v(ab)&=v(a)+v(b) &&(a,b\in K^\times),\\
v(a+b)&\ge\min\bigl\{v(a),v(b)\bigr\} &&(a,b\in K^\times,\ a+b\ne0).
\end{aligned}
\]
这样的 \(v\) 称为一个\term[lisan-fuzhi]{离散赋值}[discrete valuation]；相应子环
\(V=\{0\}\cup\Bigl\{\,a\in K^\times\mid v(a)\ge0\,\Bigr\}\) 称为\term[lisan-fuzhi-huan]{离散赋值环}[discrete valuation ring]，简称 DVR。
\end{definition}

由赋值条件，\(V\) 确实对加法及乘法封闭；其单位恰为赋值为零的元素。若选 \(\pi\) 满足 \(v(\pi)=1\)，每个非零元素为 \(\pi^r u\)，并且非单位组成 \((\pi)\)。因此它是非域的局部整环，所有非零理想均为 \((\pi^r)\)，证明与上一段的最小赋值论证相同。\(\mathbb Z_{(p)}\) 及 \(k[x]_{(x)}\) 分别对应有理数的 \(p\) 指数及有理函数在 \(x=0\) 处的零点阶数。\chapref{b3:ch:10}将证明 DVR 的其他内在刻画；当前只使用这些直接算例。

\ProseParagraphBreak
另一类局部环来自有限维代数。例如，对整数 \(r\ge1\)，取
\[
B=k[\epsilon]/(\epsilon^r),\qquad
C=k[x,y]/(x,y)^2.
\]
在 \(B\) 中，一个元素为单位当且仅当常数项非零：把它写成 \(a(1+n)\)，其中 \(a\in k^\times\)、\(n^r=0\)，逆为
\(a^{-1}\Bigl(1-n+\cdots+(-n)^{r-1}\Bigr)\)。在 \(C\) 中同理，唯一极大理想为 \((\bar x,\bar y)\)，其平方为零。这些环作为 \(k\) 向量空间有限维，故理想降链必稳定；它们是 Artin 局部环，正式结构将在下一章证明。与 DVR 不同，它们的极大理想幂会变成零，而且 \(C\) 的极大理想需要两个生成元。

\subsection{局部化、Hensel 化与完备化的比较}
\label{b3:subsec:0205:04}

局部化的任务已经明确：允许一个乘法集中的元素作分母。例如
\(k[x]_{(x)}\) 中的 \(1/(1-x)\) 是一个分式，没有因此加入所有形式幂级数。

\ProseParagraphBreak
完备化则要求同时保存模掉 \(\mathfrak m,\mathfrak m^2,\ldots\) 的相容近似，其基本对象是
\[
\widehat R=\varprojlim_n R/\mathfrak m^n.
\]
对 \(k[x]_{(x)}\)，每个商同构于 \(k[x]/(x^n)\)，相容的截断多项式逐项确定一个形式幂级数，所以其完备化为 \(k[[x]]\)。这项识别只用系数的相容性；完备化何时保持正合性、原环何时嵌入其中，将由\chapref{b3:ch:13}在明确的有限性条件下证明。

\ProseParagraphBreak
Hensel 化处理的是剩余域中的单根或互素因子能否提升到局部环的问题。在特征不为 \(2\) 的域上，方程 \(T^2-(1+x)=0\) 模掉 \(x\) 后在 \(T=1\) 处有单根；但它在 \(k[x]_{(x)}\) 中没有根，因为有理函数的平方在不可约多项式 \(1+x\) 处具有偶数阶，而 \(1+x\) 的阶为一。Hensel 化会按其泛性质补足这类提升。\chapref{b3:ch:23}才给出 Hensel 环与 Hensel 化的完整定义、构造及性质；本段只是区分三种操作的目的，不把尚未建立的提升性质用于本章证明。
